\section{Classification as a custody problem}
\label{sec:14.5}
Classification does three things to a refusal, and this section collects them. It keeps the reason from the refuser: the case description is the only channel, and Minab stays coded as it was (position 3). It keeps the record from the public: an investigation officials call complete and not released, a System Oversight section withheld in full, a review ratio no outsider can read. And it keeps the reviewer from being adverse: a cleared committee's lever is a budget line, an inspector general sits inside, and the people a deployment is aimed at can't be cleared into their own case.\footnote{A cleared legislative committee is adverse, but its lever is a budget line. An inspector general sits inside the operator. The ICRC, under its confidentiality model, reaches the belligerent and publishes nothing. A Geneva protecting power is built for a neighbouring job and is rarely designated.} Three kinds of secrecy produce these effects. Essential secrecy, where the operation fails if the enemy can read it. Accidental secrecy, where a default classifies what nobody decided to keep. And evasive secrecy, where review is moved from authorization-before to explanation-after by classifying the evidence it would need. My instruments answer them differently, and this section says which answers which.

Essential secrecy is answered by letting evidence in, not out. Classification stops information leaving, and bringing it in is routine, so the import condition, deconfliction coordinates and open-source imagery brought into the enclave as a condition of licensing, puts evidence the operator didn't compose where the system can reach it without anything leaving. A refuser inside the compartment works as the cleared bearer of \S\ref{sec:11.5} does, and the operation's secrets stay where they were.

Accidental secrecy is answered by the record. The review ratio and the committed log are lodged with an adverse party when the act occurs and published after the campaign, and class-level notice publishes the classes and counts on the same lag, so what nobody decided to keep secret comes out on a schedule nobody has to decide.

Evasive secrecy is the hard case, because there the classification is the act. It is answered by an adverse party with standing: counsel for the people a deployment affects, cleared and appearing as of right; the procedure the Classified Information Procedures Act already supplies in criminal cases (\S\ref{sec:14.7}); and the trend measure, which shows whether what dissent costs is rising while the risk stays flat.
